% ═══════════════════════════════════════════════════════════════════════════
% LEHRERHINWEISE: Elektromagnetische Induktion (Kl. 9, DS 09b)
% Kurz-Lessonplan + Lösungen + Differenzierung
% ═══════════════════════════════════════════════════════════════════════════

\documentclass[a4paper,11pt]{article}

\usepackage[utf8]{inputenc}
\usepackage[T1]{fontenc}
\usepackage[ngerman]{babel}
\usepackage{helvet}
\renewcommand{\familydefault}{\sfdefault}

\usepackage[margin=20mm]{geometry}
\usepackage{xcolor}
\usepackage{tabularx}
\usepackage{booktabs}
\usepackage{enumitem}
\usepackage{tcolorbox}
\usepackage{hyperref}
\usepackage{amssymb}

\definecolor{physik}{HTML}{2c5aa0}
\definecolor{hinweis}{HTML}{b35c00}
\definecolor{erfolg}{HTML}{2a7a4b}
\definecolor{gefahr}{HTML}{c62828}

\tcbuselibrary{skins,breakable}

\newtcolorbox{infobox}[1][]{
    colback=physik!10,
    colframe=physik,
    fonttitle=\bfseries,
    title=#1
}

\newtcolorbox{warnbox}[1][]{
    colback=hinweis!10,
    colframe=hinweis,
    fonttitle=\bfseries,
    title=#1
}

\newtcolorbox{successbox}[1][]{
    colback=erfolg!10,
    colframe=erfolg,
    fonttitle=\bfseries,
    title=#1
}

\pagestyle{empty}

\begin{document}

{\Large\bfseries\color{physik}Lehrerhinweise: Elektromagnetische Induktion}

\vspace{2mm}
{\small Klasse 9 | Einzelstunde 40' | DS 09b}

\vspace{5mm}

% ═══════════════════════════════════════════════════════════════════════════
\section*{Stundenverlauf (40 Minuten)}
% ═══════════════════════════════════════════════════════════════════════════

\begin{tabularx}{\textwidth}{|c|c|X|X|}
\hline
\textbf{Phase} & \textbf{Zeit} & \textbf{Inhalt} & \textbf{Material/Methode} \\
\hline
Einstieg & 3' & Rückblick Motor, Leitfrage: ``Geht es auch umgekehrt?'' & Plenum \\
\hline
DE & 3' & Induktion zeigen: Magnet in Spule, Voltmeter & Demo-Material \\
\hline
Erarbeitung & 12' & LP-09b durcharbeiten (Simulation!) & Tablets/Smartphones, QR-Code \\
\hline
Sicherung & 4' & Kurze Besprechung: Was ist Induktion? & Plenum \\
\hline
\textbf{Leistung} & \textbf{16'} & \textbf{MT-09b: Induktion (25 BE)} & Tablets, QR-Code \\
\hline
Abschluss & 2' & Ergebnis notieren, Ausblick Generator/Trafo & --- \\
\hline
\end{tabularx}

\vspace{5mm}

% ═══════════════════════════════════════════════════════════════════════════
\section*{Demonstrationsexperiment (Pflicht laut Rahmenplan)}
% ═══════════════════════════════════════════════════════════════════════════

\begin{infobox}[DE: Induktionsspannung erzeugen]
\textbf{Aufbau:}
\begin{itemize}[noitemsep]
    \item Spule (ca. 500--1000 Windungen)
    \item Stabmagnet (Neodym oder Ferrit)
    \item Empfindliches Voltmeter (analog oder digital)
\end{itemize}

\textbf{Durchführung:}
\begin{enumerate}[noitemsep]
    \item Magnet stillhalten $\rightarrow$ keine Spannung
    \item Magnet langsam durch Spule bewegen $\rightarrow$ kleine Spannung
    \item Magnet schnell bewegen $\rightarrow$ größere Spannung
    \item Bewegungsrichtung umkehren $\rightarrow$ Vorzeichen wechselt
    \item Magnet umdrehen $\rightarrow$ Vorzeichen wechselt
\end{enumerate}

\textbf{Erweiterung (wenn Zeit):}
\begin{itemize}[noitemsep]
    \item Faraday-Lampe (Schütteltaschenlampe) zeigen
    \item Motor als Generator nutzen: Rotor von Hand drehen
\end{itemize}
\end{infobox}

\vspace{5mm}

% ═══════════════════════════════════════════════════════════════════════════
\section*{Alternative Simulationen / Externe Ressourcen}
% ═══════════════════════════════════════════════════════════════════════════

\begin{warnbox}[Empfohlene externe Ressourcen]
\textbf{LEIFI Physik} (sehr empfohlen!)
\begin{itemize}[noitemsep]
    \item \url{https://www.leifiphysik.de/elektrizitaetslehre/elektromagnetische-induktion}
    \item Animationen zu Induktion und Generator
    \item Aufgaben mit Lösungen
\end{itemize}

\textbf{PhET Colorado}
\begin{itemize}[noitemsep]
    \item \url{https://phet.colorado.edu/de/simulations/faradays-law}
    \item ``Faraday's Law'' -- interaktive Induktions-Simulation
    \item Zeigt Magnet, Spule, Glühlampe
\end{itemize}

\textbf{YouTube (zur Vorbereitung)}
\begin{itemize}[noitemsep]
    \item SimpleClub: ``Elektromagnetische Induktion einfach erklärt''
    \item musstewissen Physik: ``Was ist Induktion?''
\end{itemize}

\textbf{Eigene Simulation}
\begin{itemize}[noitemsep]
    \item \texttt{SIM-09b-induktion.html} -- interaktiv, zeigt Spannung bei Bewegung
\end{itemize}
\end{warnbox}

\vspace{5mm}

% ═══════════════════════════════════════════════════════════════════════════
\section*{Lösungen LB-09b}
% ═══════════════════════════════════════════════════════════════════════════

\begin{successbox}[Musterlösungen Lernblatt]

\textbf{Kasten 1: Was ist Induktion?}
\begin{itemize}[noitemsep]
    \item Magneten, Spule, Spannung
    \item Änderung/Bewegung
    \item keine
    \item Faraday
\end{itemize}

\textbf{Kasten 2: Einflussfaktoren}
\begin{enumerate}[noitemsep]
    \item schneller
    \item stärker
    \item Windungen
\end{enumerate}
Richtung, Vorzeichen

\textbf{Kasten 3: Motor und Generator}
\begin{itemize}[noitemsep]
    \item Geschwister
    \item hinein, heraus
    \item Elektrische, Bewegungs-
    \item Kraftwerk, Dynamo, Windrad
\end{itemize}

\textbf{Übung 1: Anwendungen}
\begin{itemize}[noitemsep]
    \item Fahrraddynamo: Generator
    \item E-Auto (Antrieb): Motor
    \item Windkraftanlage: Generator
    \item Faraday-Lampe: Generator
\end{itemize}

\textbf{Übung 2: Energiekette}\\
Lageenergie $\rightarrow$ Bewegungsenergie $\rightarrow$ Elektrische Energie

\textbf{Übung 3: Transfer}\\
a) Die Spannung wird etwa doppelt so groß (ca. 1~V).\\
b) Die Spannung wird null (keine Bewegung = keine Induktion).

\end{successbox}

\vspace{5mm}

% ═══════════════════════════════════════════════════════════════════════════
\section*{Lösungen MT-09b (25 BE)}
% ═══════════════════════════════════════════════════════════════════════════

\begin{successbox}[Musterlösungen Mini-Test]

\textbf{Aufgabe 1: Zuordnung (5 BE)}
\begin{itemize}[noitemsep]
    \item (a) Strom fließt hinein: Motor
    \item (b) Erzeugt elektrische Energie: Generator
    \item (c) Wird im Kraftwerk verwendet: Generator
    \item (d) Wandelt el. Energie in Bewegung um: Motor
    \item (e) Funktioniert durch Induktion: Generator
\end{itemize}

\textbf{Aufgabe 2: Wann entsteht Induktionsspannung? (4 BE)}\\
\textbf{(b)} Wenn der Magnet sich relativ zur Spule bewegt

\textbf{Aufgabe 3: Einflussfaktoren (4 BE)}\\
\textbf{(c)} Längerer Draht (bei gleicher Windungszahl) erhöht die Spannung NICHT.

\textbf{Aufgabe 4: Energieumwandlung (4 BE)}\\
\textbf{(b)} Bewegungsenergie (Wasser) $\rightarrow$ Elektrische Energie

\textbf{Aufgabe 5: Anwendung (4 BE)}\\
\textbf{(b)} Ein Magnet bewegt sich beim Schütteln durch eine Spule

\textbf{Aufgabe 6: Transfer (4 BE)}\\
\textbf{(b)} Bei schnellerer Drehung wird eine größere Spannung induziert

\end{successbox}

\vspace{5mm}

% ═══════════════════════════════════════════════════════════════════════════
\section*{Differenzierung}
% ═══════════════════════════════════════════════════════════════════════════

\begin{tabularx}{\textwidth}{|c|X|X|}
\hline
\textbf{Niveau} & \textbf{Unterstützung} & \textbf{Erweiterung} \\
\hline
$\star$ (G-Basis) &
Simulation mit Hilfestellung; Begriffe vorgeben &
Nur Aufgaben 1, 2, 4, 5 bearbeiten (17 BE) \\
\hline
$\star\star$ (Standard) &
LP vollständig durcharbeiten &
Alle Aufgaben (25 BE) \\
\hline
$\star\star\star$ (E-Erw.) &
--- &
Zusatzfrage: ``Erkläre den Aufbau eines Generators.'' \\
\hline
\end{tabularx}

\vspace{5mm}

% ═══════════════════════════════════════════════════════════════════════════
\section*{Häufige Schülerfehler}
% ═══════════════════════════════════════════════════════════════════════════

\begin{itemize}[noitemsep]
    \item \textbf{``Induktion gibt es auch ohne Bewegung'':} Nein! Ohne Änderung keine Spannung.
    \item \textbf{Verwechslung Motor/Generator:} Motor: Strom rein $\rightarrow$ Bewegung raus. Generator umgekehrt!
    \item \textbf{``Schnellere Bewegung = weniger Spannung'':} Umgekehrt! Schneller = größere Spannung.
    \item \textbf{Energieumwandlung unklar:} Generator: Bewegung $\rightarrow$ Strom. Nicht umgekehrt!
\end{itemize}

\vspace{5mm}

% ═══════════════════════════════════════════════════════════════════════════
\section*{Materialliste}
% ═══════════════════════════════════════════════════════════════════════════

\begin{itemize}[noitemsep]
    \item[$\square$] Spule (500--1000 Windungen)
    \item[$\square$] Stabmagnet
    \item[$\square$] Empfindliches Voltmeter
    \item[$\square$] Optional: Faraday-Lampe (Schütteltaschenlampe)
    \item[$\square$] QR-Code: LP-09b (Projektion oder Ausdruck)
    \item[$\square$] QR-Code: MT-09b (Projektion)
    \item[$\square$] Tablets/Smartphones für SuS
\end{itemize}

\vspace{5mm}

% ═══════════════════════════════════════════════════════════════════════════
\section*{Anschluss: Nächste Stunden (Generator, Transformator)}
% ═══════════════════════════════════════════════════════════════════════════

\begin{infobox}[Überleitung]
\textbf{Mögliche Folgestunden (wenn Zeit):}
\begin{itemize}[noitemsep]
    \item Generator: Aufbau, Wechselspannung
    \item Transformator: Spannungs-/Stromumwandlung
    \item Anwendungen: Kraftwerke, Stromnetz
\end{itemize}

\textbf{Für Klasse 9 lt. Rahmenplan ausreichend:}
\begin{itemize}[noitemsep]
    \item Qualitatives Verständnis von Induktion
    \item Vergleich Motor/Generator (Umkehrung)
    \item Energieumwandlung benennen können
\end{itemize}

Generator und Transformator werden ``in Form von Analogien'' behandelt -- keine quantitativen Berechnungen in Klasse 9!
\end{infobox}

\end{document}
